% =====================================================================
% Appendices
% Each appendix begins with \chapter{...} (the \appendix command in
% main.tex switches \chapter numbering to A, B, C, ...).
% =====================================================================

\chapter{Declaration of Use of Generative AI}
\label{app:ai-declaration}

% ---------------------------------------------------------------------
% This appendix is REQUIRED. DTU permits the use of generative AI but
% mandates that students declare distinctly where and how it was used.
% The human authors remain the sole authors; AI is not a co-author.
% Sources: DTU library guidance and DTU AI rules (see CLAUDE.md §5).
% ---------------------------------------------------------------------

\section{Statement of Authorship and Responsibility}

We, the authors of this thesis, take full and sole responsibility for the content of this report.
All decisions regarding scope, claims, results, and interpretations are ours.
The generative artificial intelligence (GAI) tools listed below were used in a limited, supporting role; they are not co-authors of this work.
In line with DTU guidance, this appendix records which tools we used, and where and how their output entered our process, so that the reader has a clear and bounded understanding of that contribution.

\section{Tools Used}

\begin{table}[!htbp]
\centering
\small
\setlength{\tabcolsep}{4pt}
\renewcommand{\arraystretch}{1.2}
\begin{tabular}{@{}p{3.4cm}p{3.6cm}p{7.0cm}@{}}
\toprule
\textbf{Tool} & \textbf{Version / Period} & \textbf{Purpose} \\
\midrule
ChatGPT (OpenAI) & GPT-5.5 (latest available during Feb--June 2026) & Feedback on our own draft prose; discussion of ideas and design trade-offs. \\
\addlinespace
Codex (OpenAI) & GPT-5.5 (latest available during Feb--June 2026) & Coding and design assistance for the software artifact. \\
\addlinespace
Claude (Anthropic) & Opus 4.8 and Sonnet 4.6 (latest available during Feb--June 2026) & Coding and design assistance; sounding board for design trade-offs; paraphrasing of our draft prose; LaTeX scaffolding. \\
\bottomrule
\end{tabular}
\caption{Generative AI tools used during the preparation of this thesis.}
\label{tab:gai-tools}
\end{table}

\section{Where and How GAI Was Used}
\label{sec:ai-how-used}

\subsection{Idea Validation and Trade-off Discussion}
\label{sec:ai-tradeoff}

Our principal use of GAI was as a sounding board to stress-test reasoning we had already formed, rather than to originate it.
A typical interaction was to ask a tool to argue against a design choice we were considering, so that we could check whether we had overlooked a counter-argument.
This applied in particular to the recurring tension between architectural uniformity and responsiveness that is consolidated in the design decision register (\cref{sec:design-decisions}).
No design decision was adopted on the basis of a tool's output; every decision, and the argument defending it, was formed, verified, and written by us.

\subsection{Writing Assistance}

We used GAI to obtain feedback on prose we had already written and to suggest alternative phrasings for clarity and concision, across the main content chapters of this report.
In particular, we used Claude to paraphrase passages we had already drafted; the report and the meaning behind it remain entirely our own.
We did not use it to generate sections from scratch; the prose is our own, and AI assistance was limited to feedback and rephrasing of text we had already drafted.
All factual claims, citations, and numbers in any AI-assisted passage were verified by us against their original sources, and the final wording is our own.

\subsection{Code Assistance}

The accompanying software artifact (the \texttt{reading-the-reader} monorepo) was developed with coding and design assistance from Codex and Claude, used for scaffolding, refactoring suggestions, design discussion, and debugging across the React and TypeScript frontend and the C\# and .NET backend.
We also used it for LaTeX scaffolding and build configuration of this report.
All generated code was reviewed, tested, and integrated by us, and we take responsibility for its correctness.

\section{What GAI Was Not Used For}

To make the boundary explicit, the following were done entirely by us, without GAI:

\begin{itemize}
    \item Formulating the problem statement and the research questions.
    \item Selecting, reading, and interpreting the cited literature (no source was cited on the basis of an AI-generated summary alone).
    \item Making the system's design and architecture decisions, and constructing the arguments that defend them (GAI served only as a sounding board to surface trade-offs, as described in \cref{sec:ai-tradeoff}).
    \item Designing the evaluation, conducting it, and interpreting the results.
    \item Drawing the conclusions and identifying the limitations and threats to validity.
\end{itemize}

% =====================================================================
% Additional appendices below (extended results, full requirements
% list, raw evaluation data, repository structure, ethics approval,
% etc.). Add new chapters as needed.
% =====================================================================

% \chapter{Full Requirements Specification}
% \label{app:requirements}

% \chapter{Repository Structure}
% \label{app:repo}

% \chapter{Raw Evaluation Data}
% \label{app:raw-data}
